% ID: SEL-D0022-OB-E02
% ESTADO: validada. S023, revisada con los criterios oficiales.
\textbf{a)} Cada tramo es polinómico. En $t=10$ el límite izquierdo
es $40$ y el valor del segundo tramo es $-20+80-20=40$.
Por tanto es continua en todo $[0,25]$. Las pendientes interiores
son $4$ y $8-2t/5$, que coinciden en $10$. Es derivable en $(0,25)$,
con derivadas laterales $4$ en $0$ y $-2$ en $25$.

\textbf{b)} El segundo tramo tiene derivada positiva antes de
$t=20$ y negativa después. En conjunto, $B$ crece en $[0,20]$
y decrece en $[20,25]$. El máximo llega en el
$\boxed{\text{año }20}$ y vale $\boxed{60\text{ miles de euros}}$.

\textbf{c)} Un segmento une $(0,0)$ con $(10,40)$.
Le sigue el arco cóncavo $60-(t-20)^2/5$, que pasa por
$(20,60)$ y termina en $(25,55)$.
\begin{center}
\setlength{\unitlength}{3mm}
\begin{picture}(31,16)(-3,-2)
\qbezier(0,0)(13,0)(26,0)\put(26.5,0){$t$}
\qbezier(0,0)(0,6.5)(0,13)\put(0,13.2){$B$}
\qbezier(0,0)(5,4)(10,8)
\qbezier(10,8)(17.5,14)(25,11)
\put(10,7.6){\makebox(0,0)[t]{$(10,40)$}}
\put(20,12.5){\makebox(0,0)[b]{$(20,60)$}}
\put(25,10.7){\makebox(0,0)[l]{$(25,55)$}}
\end{picture}
\end{center}
